\documentclass[preview]{standalone}

\usepackage{amsmath}
\usepackage{amssymb}
\usepackage{stellar}
\usepackage{definitions}
\usepackage{bettelini}

\begin{document}

\id{parametrized-surfaces}
\genpage

\section{Parametrized Surfaces}

\begin{snippet}{surface-integration-motivation}
    Let \(\Sigma\subseteq\realnumbers^3\) be a surface and let
    \(f\colon\Sigma\fromto\realnumbers\). To give meaning to
    \[
        \int_\Sigma f\,dS,
    \]
    one describes \(\Sigma\) by parameters and reduces the integral to an
    ordinary double integral on a subset of \(\realnumbers^2\).
\end{snippet}

\begin{snippetdefinition}{parametrized-surface-definition}{Parametrized surface}
    Let \(A\subseteq\realnumbers^k\) be bounded and open. A
    \(k\)-dimensional parametrization in \(\realnumbers^n\) is a map
    \[
        \sigma\colon\overline A\fromto\realnumbers^n.
    \]
    Its image \(\Sigma=\sigma(\overline A)\) is called a
    \emph{parametrized \(k\)-surface}, or a parametrized
    \(k\)-dimensional submanifold.
\end{snippetdefinition}

\begin{snippet}{surface-parametrization-three-dimensional-case}
    For a surface in \(\realnumbers^3\), one has \(k=2\), \(n=3\), and
    usually writes
    \[
        \sigma(u,v)=\bigl(x(u,v),y(u,v),z(u,v)\bigr).
    \]
\end{snippet}

\begin{snippetdefinition}{regular-surface-parametrization-definition}{Regular parametrization}
    Let
    \[
        \sigma\colon\overline A\subseteq\realnumbers^k
        \fromto\realnumbers^n.
    \]
    The parametrization is \emph{regular} if
    \(\sigma\in\mathcal C^1(\overline A)\) and
    \[
        \operatorname{rank}J_\sigma(t)=k
        \qquad\text{for every }t\in A.
    \]
    Equivalently, the vectors
    \(\partial_{t_1}\sigma(t),\ldots,\partial_{t_k}\sigma(t)\) are
    linearly independent.
\end{snippetdefinition}

\begin{snippet}{surface-jacobian-matrix}
    If \(t=(t_1,\ldots,t_k)\), then
    \[
        J_\sigma(t)=
        \begin{bmatrix}
            \partial_{t_1}\sigma(t)&\cdots&\partial_{t_k}\sigma(t)
        \end{bmatrix}.
    \]
    In the case \(k=2\), \(n=3\),
    \[
        J_\sigma(u,v)=
        \begin{bmatrix}
            x_u&x_v\\
            y_u&y_v\\
            z_u&z_v
        \end{bmatrix}.
    \]
\end{snippet}

\begin{snippetproposition}{surface-regularity-geometric-interpretation}{Geometric interpretation of regularity}
    Let
    \(\sigma\colon\overline A\subseteq\realnumbers^2
    \fromto\realnumbers^3\) be regular and let
    \(p=\sigma(u_0,v_0)\). Then
    \[
        T_p\Sigma=
        \operatorname{span}\{\sigma_u(u_0,v_0),\sigma_v(u_0,v_0)\}.
    \]
\end{snippetproposition}

\begin{snippetproof}{surface-regularity-geometric-interpretation-proof}{surface-regularity-geometric-interpretation}{Geometric interpretation of regularity}
    The coordinate curves
    \[
        \gamma(u)=\sigma(u,v_0),
        \qquad
        \varphi(v)=\sigma(u_0,v)
    \]
    lie on the surface and pass through \(p\). Their tangent vectors at
    \(p\) are
    \[
        \gamma'(u_0)=\sigma_u(u_0,v_0),
        \qquad
        \varphi'(v_0)=\sigma_v(u_0,v_0).
    \]
    Regularity makes these vectors linearly independent, so they span the
    tangent plane.
\end{snippetproof}

\section{Simple Surfaces and Local Graphs}

\begin{snippetdefinition}{simple-surface-parametrization-definition}{Simple parametrization}
    A parametrization
    \(\sigma\colon\overline A\fromto\realnumbers^n\) is \emph{simple} if
    \[
        \sigma(t)=\sigma(s),\quad t\neq s
        \quad\Longrightarrow\quad t,s\in\partial A.
    \]
    In particular, \(\sigma\) is injective on \(A\).
\end{snippetdefinition}

\begin{snippet}{simple-parametrization-interpretation}
    Simplicity prevents the interior of the parameter domain from covering
    the same part of a surface more than once. Otherwise a parameter-domain
    integral would count that area repeatedly.
\end{snippet}

\begin{snippetexample}{graph-surface-regular-parametrization-example}{Graph of a function}
    Let \(A\subseteq\realnumbers^2\) be open and
    \(f\in\mathcal C^1(A)\). The map
    \[
        \sigma(x,y)=\bigl(x,y,f(x,y)\bigr)
    \]
    parametrizes the graph of \(f\), and
    \[
        J_\sigma(x,y)=
        \begin{bmatrix}
            1&0\\
            0&1\\
            f_x&f_y
        \end{bmatrix}.
    \]
    Its columns are linearly independent everywhere, hence the graph
    parametrization is regular.
\end{snippetexample}

\begin{snippetproposition}{regular-surface-local-graph-proposition}{Local graph representation}
    Let
    \(\sigma\colon\overline A\subseteq\realnumbers^k
    \fromto\realnumbers^n\) be regular at \(t_0\in A\). Then, after
    possibly reordering the coordinates of \(\realnumbers^n\), the image of
    \(\sigma\) is locally the graph of a function of \(k\) variables.
\end{snippetproposition}

\begin{snippetproof}{regular-surface-local-graph-proposition-proof}{regular-surface-local-graph-proposition}{Local graph representation}
    Since \(J_\sigma(t_0)\) has rank \(k\), some \(k\times k\) minor is
    nonzero. Projecting onto the corresponding \(k\) target coordinates
    gives a map with invertible derivative at \(t_0\). The inverse function
    theorem provides a local inverse for those coordinates; substituting it
    into the remaining components of \(\sigma\) gives the desired graph.

    For \(\sigma(u,v)=(x(u,v),y(u,v),z(u,v))\), if
    \(\partial(x,y)/\partial(u,v)\neq0\), then locally
    \(u=u(x,y)\), \(v=v(x,y)\), and the surface is
    \[
        (x,y)\longmapsto z\bigl(u(x,y),v(x,y)\bigr).
    \]
\end{snippetproof}

\begin{snippetdefinition}{surface-coordinate-curves-definition}{Coordinate curves}
    Let
    \(\sigma\colon\overline A\subseteq\realnumbers^k
    \fromto\realnumbers^n\). A \emph{coordinate curve} is obtained by
    fixing \(k-1\) parameters and varying the remaining one.

    If \(k=2\), the coordinate curves through
    \(\sigma(u_0,v_0)\) are
    \[
        u\longmapsto\sigma(u,v_0),
        \qquad
        v\longmapsto\sigma(u_0,v),
    \]
    with velocities \(\sigma_u(u_0,v_0)\) and
    \(\sigma_v(u_0,v_0)\), respectively.
\end{snippetdefinition}

\end{document}
